% HD-F1-T23. GPT-6 Astra Pro; /root/transport; 22 November 2026 UTC.
\section[HD-F1-T23: complete alternating-socle exclusion over seven fields]{HD-F1-T23: the actual central-cover maps complete the alternating-socle exclusion over seven fields}
\label{sec:f1-alternating-socle-completion}

For the seven fields
\[
 K=\mathbf Q(\sqrt{-d}),\qquad
 d\in\{2,7,11,19,43,67,163\},
 \tag{F1T23.1}
\]
the remaining natural sixfold source has no almost
simple projective coefficient image with alternating
socle \(A_n\), for any \(n\ge5\).
Over \(\mathbf Q(i)\), the same exact central-cover
argument and T21 leave only an ordinary layer
\(L\simeq A_n\), \(n\ge7\), in the alternating
alternative. This last sentence is a reduction of its
precise scope, not an exclusion of that Gaussian case.

\paragraph{1. The actual perfect central layer.}
Let \(G\subset GL_K(V)\) be the original faithful
absolutely irreducible coefficient group, and let
\(Z=Z(G)=G\cap K^*I\) be its scalar center.
Let \(P\) be the inverse image of the socle in
\(G/Z\), and use its actual quotient map
\(\pi:P\to A_n\). Perfectness of \(A_n\)
gives \(\pi(P')=A_n\); hence
\[
 P=ZL,\qquad L=P',\qquad
 L=[P,P]=[ZL,ZL]=[L,L].
 \tag{F1T23.2}
\]
Thus \(L\) is an actual perfect central extension,
not a newly chosen covering group. Its kernel is
\(C=L\cap Z\). Since \(A_n\) is centerless,
\(Z(L)=C\): a central element has central image in
\(A_n\), and the reverse containment is immediate.
In (F1T23.1), \(C\subset\{1,-1\}\); over
\(\mathbf Q(i)\), \(C\subset\mu_4\).

\paragraph{2. The explicit universal-cover map onto that layer.}
Here is the full map, including its independence from
choices of lifts. Choose a free presentation
\(F\twoheadrightarrow A_n\), with kernel \(R\),
and put
\[
 E=F/[F,R],\qquad U=E',\qquad
 p:U\longrightarrow A_n,\qquad M=\ker p.
 \tag{F1T23.3}
\]
The kernel of \(E\to A_n\) is central.
As in (F1T23.2), perfectness of the quotient gives
\(E=(R/[F,R])U\), \(U=[U,U]\), and surjectivity
of \(p\). Choose arbitrary lifts in \(L\) of the
free generators. The resulting homomorphism
\(F\to L\) sends \(R\) into \(C\), so kills
\([F,R]\) and descends to \(E\to L\).
Its restriction is the actual homomorphism
\[
 \beta:U\longrightarrow L\hookrightarrow G,
       \qquad\pi\beta=p.
 \tag{F1T23.4}
\]
The image of \(E\) together with \(C\) generates
\(L\), because it maps onto \(A_n\). Taking
commutators and using centrality of \(C\) and
perfectness of \(L\) proves
\(\beta(U)=L\). Different lift choices multiply
each free-generator image by a central element, so
give identical maps on \(F'\), hence on \(U\).
More generally, any two homomorphisms from \(U\)
to a central extension lifting \(p\) have a pointwise
central quotient. That quotient is a homomorphism
from the perfect \(U\) to an abelian group, so is
trivial. This proves the universal property and
identifies (F1T23.3) with the universal central cover.
The construction and these facts are also the full
proofs of Aschbacher (33.3)--(33.8), cited below.

Since \(\pi\beta=p\), the kernel
\(N=\ker\beta\) lies in \(M\). Surjectivity
gives the exact quotient maps
\[
       U/N\xrightarrow{\sim}L,\qquad
       M/N\xrightarrow{\sim}C=Z(L),
       \qquad uN\longmapsto\beta(u).
 \tag{F1T23.5}
\]
These equalities identify the center of the original
matrix layer with a specific quotient of the multiplier.

\paragraph{3. The published multiplier theorem and the exhaustive permitted quotients.}
Use M.~Aschbacher, \emph{Finite Group Theory},
second edition, Cambridge University Press (2000),
\S33, (33.15)(1), printed pp.170--171:
\[
 M(A_n)\simeq
 \begin{cases}
 C_6,&n=6,7,\\
 C_2,&n=5\text{ or }n\ge8.
 \end{cases}
 \tag{F1T23.6}
\]
The complete indexed primary text of \S33,
including the universal-cover proofs and the supplied
multiplier proof through p.174, was read. Its
identification of the extra factor three for \(n=6,7\)
is a stated classical theorem; that proof is explicitly
omitted in the book, and is not claimed to have been
reconstructed here. The needed center restriction
already rules out that factor. The precise access
scope is recorded in the source note.

For \(n\ge8\), (F1T23.5)--(F1T23.6) leave
\(N=M\), giving \(L\simeq A_n\), or \(N=1\),
giving the standard double cover \(L\simeq2.A_n\).
For \(n=7\), a cyclic group of order six has one
subgroup of each divisor order. The constraint
\(C\subset\{1,-1\}\) leaves exactly
\[
 \begin{array}{c|c|c}
 M&N&L=U/N\\\hline
 C_6&C_6&A_7\\
 C_6&C_3&2.A_7.
 \end{array}
 \tag{F1T23.7}
\]
In particular a surviving central element of order
three would have to act as a scalar cube root of
unity in \(K\); faithfulness forbids killing it
inside the original matrix image. Such a scalar
does not exist in the fields (F1T23.1).

The identification with the standard double cover is
also an exact quotient identification. That standard
perfect cover maps to \(A_n\) with kernel of order
two, so universality gives an epimorphism from \(U\)
onto it. Its kernel is the unique subgroup of \(M\)
of index two. This is the same subgroup \(N\) in
(F1T23.5); the quotient maps therefore give an
isomorphism over \(A_n\) from the standard cover
to the original \(L\). The standard perfect double
cover and its subgroup maps are the published ones
already checked in T21. This argument proves the
required identification rather than assigning the
name \(2.A_n\) merely from the order of the center.

T22 excludes the ordinary outcome for every \(n\ge7\)
over (F1T23.1). T21 excludes the double-cover outcome
for every \(n\ge7\) and every remaining coefficient
field. Over \(\mathbf Q(i)\), the same list applies:
the center has order dividing four and is a quotient
of \(C_2\) or \(C_6\), so its order again divides
two. T21 then leaves only the ordinary Gaussian layer.

\paragraph{4. The small alternating socles use the already proved all-cover case.}
T14 excludes every actual almost simple projective
socle \(PSL_2(q)\), \(q\ge4\), including all
central lifts and all outer inertia. Its proof uses
the original layer and requires no multiplier bound.
The identifications for \(A_5\) and \(A_6\) can
be specified as follows.

The fractional-linear action
\[
 PSL_2(4)\longrightarrow\operatorname{Perm}
             (\mathbf P^1(\mathbf F_4))=S_5,
 \qquad z\longmapsto\frac{az+b}{cz+d},
 \tag{F1T23.8}
\]
is faithful. Its order is sixty; its image is even
because this simple group has no quotient of order
two. Thus its image is exactly \(A_5\), giving
the first isomorphism with its inverse.

For the second, the primary ATLAS \(A_6\) record
gives generators \(a=(12)(34)\),
\(b=(1235)(46)\) and the presentation
\(a^2=b^4=(ab)^5=(ab^2)^5=1\).
Let \(\alpha\in\mathbf F_9\) satisfy
\(\alpha^2=\alpha+1\), so \(\alpha^4=-1\),
and take the following actual ATLAS matrices:
\[
 A=\begin{pmatrix}\alpha^6&0\\\alpha^3&\alpha^2\end{pmatrix},
 \qquad
 B=\begin{pmatrix}0&1\\-1&\alpha^2\end{pmatrix}
                  \quad\text{in }SL_2(9).
 \tag{F1T23.9}
\]
They satisfy \(A^2=B^4=-I\).
Their product traces are
\(\operatorname{tr}(AB)=2\alpha\) and
\(\operatorname{tr}(AB^2)=\alpha+2\);
both solve \(t^2+t-1=0\). For a determinant-one
matrix with such a trace, the relation
\(X^2-tX+I=0\) gives \(X^5=I\) by direct
substitution. Thus
\[
        A_6\longrightarrow PSL_2(9),
                   \qquad a\longmapsto[A],\quad
                          b\longmapsto[B]
 \tag{F1T23.10}
\]
is a well-defined nontrivial homomorphism. Simplicity
of \(A_6\) makes it injective, and both groups
have order360, so it is an isomorphism. The stated
presentation and matrices are separately linked in
the source record. In the GAP matrices the primitive
element is denoted \(Z(9)\); (F1T23.9) fixes the
field presentation explicitly and verifies the
relations directly.

Conjugation by either isomorphism \(\phi\)
transports an actual projective automorphism by
\(x\mapsto\phi x\phi^{-1}\). This identifies
the original almost simple image with one of the
literal \(PSL_2\) cases of T14 while keeping its
original inverse image \(P\), commutator \(L\),
and coefficient matrices. T14 excludes both small
socles. Together with paragraph3 this completes the
seven-field alternating-socle exclusion for all
\(n\ge5\).

\paragraph{Original geometric maps and source scope.}
The only new group maps are the free-group lift,
the universal-cover surjection and its central
quotients, and the displayed small-socle
isomorphisms. All are compared inside the original
coefficient image before T14, T21 or T22 is used.
The original connected curve cover, primitive
rational projector, sixfold corner, integral-cleared
Jacobian image and restricted theta polarization
retain their identity maps. Their Poincare
biextension, cup form, de Rham and Hodge
realizations, period derivative, Hilbert and
Picard maps, cycle and Grothendieck classes, and
Adams operations retain multiplier one. The
result removes source alternatives; it supplies
no additional detector component.
